\documentclass[11pt,a4paper]{article}

\usepackage{amsmath}
\usepackage{mathtools}
\usepackage{fontspec}
\usepackage{unicode-math}
\setmathfont{KpMath-Regular.otf}
\setmathfont{KpMath-Bold.otf}[version=bold]
\usepackage{microtype}
\usepackage[dvipsnames]{xcolor}
\usepackage[margin=25mm, top=30mm, bottom=30mm]{geometry}

\usepackage{fancyhdr}
\pagestyle{fancy}
\fancyhf{}
\fancyhead[L]{\small\itshape Multiresolution Representation of the Background}
\fancyhead[R]{\small\itshape Nestor (2026)}
\fancyfoot[C]{\thepage}
\renewcommand{\headrulewidth}{0.4pt}
\renewcommand{\footrulewidth}{0pt}

\setcounter{secnumdepth}{3}
\usepackage{booktabs}
\usepackage{array}
\usepackage[font=small, labelfont=bf, labelsep=period]{caption}
\usepackage[authoryear,round]{natbib}
\usepackage{setspace}
\setstretch{1.15}
\setlength{\parskip}{0.4em}
\usepackage[colorlinks=true, linkcolor=blue!60!black,
            citecolor=blue!60!black, urlcolor=blue!60!black]{hyperref}

\newcommand{\bA}{\mathbf{A}}
\newcommand{\bI}{\mathbf{I}}
\newcommand{\bb}{\mathbf{b}}
\newcommand{\avg}[1]{\langle #1 \rangle}

\title{\bfseries Multiresolution Representation of the Background Medium\\[4pt]
\large Haar MRA, Scale Refinement Inside a Krylov Solve,\\
and What May Legitimately Be Averaged}
\author{}
\date{26 July 2026}

\begin{document}
\maketitle

\begin{abstract}
\noindent
The proposal is to represent the background medium in a Haar multiresolution
analysis, relaxing it to long scales during early GMRES iterates and
reintroducing finer scales later.  This note assesses it.  Three findings.
First, the \emph{direction} conflicts with inexact-Krylov theory, which
licenses accurate matrix--vector products \emph{early} and relaxed ones
\emph{late} --- the reverse; the idea belongs one level up, as nested
iteration between solves rather than scale refinement within one Krylov space.
Second, it has been done: \citet{hustedt2004} built exactly this for
frequency-domain seismic modelling in 2004, using a wavelet multiresolution to
combine coarse- and fine-grid wavefields.  Third, and most useful, a Haar
average \emph{is} the correct coarsening --- but of the entries of the system
matrix $\bA$, not of $(\alpha,\beta,\rho)$.  Measured here: averaging $\bA$
converges as $\mathcal{O}\bigl((H/\lambda)^2\bigr)$, while averaging velocities
is $21\times$ worse at the same wavelength and diverges as the wavelength
shortens.
\end{abstract}

\tableofcontents

% =============================================================================
\section{Haar multiresolution analysis}

A multiresolution analysis is a nested ladder of closed subspaces
$V_j \subset V_{j+1}$ of $L^2$ with $\overline{\bigcup_j V_j} = L^2$,
$\bigcap_j V_j = \{0\}$, the dyadic scaling property
$f(x)\in V_j \Leftrightarrow f(2x)\in V_{j+1}$, and an orthonormal basis of
$V_0$ generated by the integer translates of a scaling function $\varphi$
\citep{mallat1989,daubechies1992}.  The Haar case is
$\varphi = \chi_{[0,1)}$, so $V_j$ consists of functions piecewise constant on
dyadic intervals of length $2^{-j}$, with wavelet
$\psi = \chi_{[0,1/2)} - \chi_{[1/2,1)}$ spanning
$W_j = V_{j+1}\ominus V_j$ and two-scale relations
\begin{equation}\label{eq:refine}
  \varphi(x) = \varphi(2x) + \varphi(2x-1),
  \qquad
  \psi(x) = \varphi(2x) - \varphi(2x-1).
\end{equation}
Decomposition is running averages and differences: orthonormal, perfect
reconstruction, $\mathcal{O}(N)$.

Two properties matter here.  Haar has exactly \emph{one} vanishing moment, so
it represents piecewise-constant data exactly and smooth data poorly --- an
awkward fit for a background that is by construction the smooth depth average.
And the refinement relation~\eqref{eq:refine} is already familiar in this
project: the subdivision identity used by the face-contact arbiter, splitting a
cube into eight half-cubes with
$\mathrm{table}(n) = \tfrac18\sum_m \mathrm{mult}(m)\,\mathrm{table}(m)$, is
exactly~\eqref{eq:refine} in three dimensions.  The per-voxel contrast field is
likewise a tensor-product Haar $V_j$ at a single level.

% =============================================================================
\section{Where the refinement can go}

\subsection{Not inside one Krylov space}

The natural reading of the proposal --- coarse operator early, fine operator
late --- runs against the established theory of inexact Krylov methods.
\citet{bouras2005} and \citet{simoncini2003} show that GMRES tolerates a
perturbation $\Delta\bA_k$ in the matrix--vector product at step $k$ provided
\begin{equation}\label{eq:relax}
  \lVert \Delta\bA_k\rVert \;\lesssim\; \frac{\eta}{\lVert r_{k-1}\rVert}\,
  \varepsilon ,
\end{equation}
so the admissible inaccuracy \emph{grows} as the residual falls: be accurate
early, sloppy late.  The mechanism is that the perturbation enters the residual
gap weighted by the Arnoldi solution coefficient $|y_k|$, and those decay as
the iteration converges --- an error committed in the first few basis vectors
propagates into everything built on them.  Reported gains are real but modest,
$1.5$--$2.3\times$ on Laplace and $2.7$--$3.3\times$ on Stokes problems
\citep{vandeneshof2004}.

There is a second objection.  If the operator itself changes between iterates,
the space being built is no longer a Krylov space for a fixed operator and the
minimal-residual property is lost, unless the variation is handled explicitly
as a flexible \citep{saad1993} or inexact method.

\subsection{But nested iteration is exactly right}

The idea is sound one level up.  Converge on the coarse medium, then use that
solution as the \emph{initial guess} for the finer one: nested iteration, or
cascadic multigrid \citep{bornemann1996}, in which no coarse-grid correction is
propagated back.  This is between solves rather than within one, so neither
objection applies.  The same logic at the inversion level is multiscale
full-waveform inversion \citep{bunks1995}, where low frequencies are fitted
before high --- among the most influential ideas in the field.

% =============================================================================
\section{Prior art}

The specific construction has been built.  \citet{hustedt2004} present a
multi-level direct--iterative solver for frequency-domain seismic modelling in
which coarse- and fine-grid wavefields are combined either by space
interpolation (DISS) or through the multiresolution scaling property of an
orthogonal wavelet expansion (DIWS).  They report that because the wavelet
transform accounts for grid interactions, phase-shift artefacts are greatly
reduced and significantly fewer iterations are required than with space
interpolation.

One difference matters: they multiresolve the \emph{wavefield and grid}, not
the \emph{medium}.  Coarsening the medium is the harder problem, and it is
where the real constraint lies.

% =============================================================================
\section{What may legitimately be averaged}

\subsection{The question}

A Haar projection of the background replaces fine structure by block averages.
Of \emph{what}?  Averaging $(\alpha,\beta,\rho)$ over a block is the obvious
choice and is wrong: for a finely layered medium the correct long-wavelength
limit is Backus averaging \citep{backus1962}, in which compliances average
harmonically and the effective medium is transversely isotropic even when every
constituent layer is isotropic.  The general non-periodic statement is
Capdeville's homogenisation \citep{capdeville2010,capdeville2018}, where the
scale separation is a spatial low-pass filter whose cutoff is set by the
minimum wavelength and a precision factor, and the resulting smooth effective
medium is what the waves actually see.

\subsection{The right variables are the entries of $\bA$}

The layered problem marches as $\partial_z \bb = \bA(z)\,\bb$, so a stack of
thin layers has response $\prod_i \exp(\bA_i h_i)$.  As $h_i \to 0$ the Lie
product formula gives
\begin{equation}\label{eq:lie}
  \prod_i \exp(\bA_i h_i) \;\longrightarrow\;
  \exp\Bigl(\textstyle\sum_i \bA_i h_i\Bigr) = \exp\bigl(\avg{\bA} H\bigr),
\end{equation}
so the correct leading-order coarsening is the \emph{arithmetic} average of the
system matrix.  This is not a coincidence: $\bA$ is linear in
\[
  \gamma = \frac{\lambda}{\lambda+2\mu},\quad
  a = \frac{1}{\lambda+2\mu},\quad
  b = \frac{1}{\mu},\quad
  \zeta = \frac{4\mu(\lambda+\mu)}{\lambda+2\mu},\quad
  \chi = \frac{2\mu\lambda}{\lambda+2\mu},\quad \mu,\ \rho,
\]
and these are precisely the combinations Backus averages arithmetically:
$C_{33} = 1/\avg{a}$, $C_{44} = 1/\avg{b}$, $C_{66} = \avg{\mu}$,
$C_{13} = \avg{\gamma}/\avg{a}$, $\rho_{\text{eff}} = \avg{\rho}$.

\emph{Haar averaging is therefore correct --- in the variables of $\bA$, and
wrong in $(\alpha,\beta,\rho)$.}

\subsection{Measured}

A fixed realisation of a strongly contrasted fine layered medium ($256$ layers,
$\alpha = 4.0\pm1.2$~km/s, $\rho = 2.6\pm0.4$~g/cm$^3$, total thickness
$H = 1$~km) is coarsened to a single homogeneous layer two ways and compared
against the exact ordered product.  The long-wavelength limit is approached by
lowering the frequency, so the medium never changes and only $H/\lambda$ varies.

\begin{center}
\begin{tabular}{rrcc}
\toprule
$f$ (Hz) & $H/\lambda_P$ & velocity average & $\avg{\bA}$ average \\
\midrule
$0.25$ & $0.062$ & $9.95\times10^{-3}$ & $\mathbf{4.72\times10^{-4}}$ \\
$0.50$ & $0.125$ & $2.94\times10^{-2}$ & $1.78\times10^{-3}$ \\
$1.00$ & $0.249$ & $1.22\times10^{-1}$ & $9.34\times10^{-3}$ \\
$2.00$ & $0.499$ & $3.54\times10^{-1}$ & $3.01\times10^{-2}$ \\
$4.00$ & $0.998$ & $5.71\times10^{-1}$ & $1.58\times10^{-1}$ \\
$8.00$ & $1.996$ & $8.83\times10^{-1}$ & $1.29\times10^{-1}$ \\
\bottomrule
\end{tabular}
\end{center}

The $\avg{\bA}$ average converges as $\mathcal{O}\bigl((H/\lambda)^2\bigr)$ ---
the error roughly quadruples as $H/\lambda$ doubles --- while the velocity
average is $21\times$ worse at the longest wavelength and degrades to complete
failure within one wavelength.

\paragraph{What the residual is.}  Equation~\eqref{eq:lie} is first-order
Magnus; the leading correction comes from the non-commutativity of the
$\bA_i$, and it is exactly what Backus supplies in closed form as the
non-linear term in
$C_{11} = \avg{\zeta} + \avg{\gamma}^2/\avg{a}$, which a plain arithmetic
average of $\bA$ cannot produce.  In the realisation above that term is
$13\%$ of $C_{11}$, and it is what induces the anisotropy
$C_{11}/C_{33} = 1.126$, $C_{66}/C_{44} = 1.148$.  A multiresolution scheme
that wants better than second order must carry it.

\emph{Caveat.}  The Backus constants quoted here are evaluated from the closed
form; a transversely isotropic propagator was \emph{not} independently
validated in this exercise, so the anisotropy figures are analytic rather than
measured.

% =============================================================================
\section{The constraint from the coupling operator}

There is a further, project-specific limit.  The inter-scatterer coupling is
extraordinarily sensitive to the reference: holding geometry fixed and varying
only whether the crust reference carries its depth gradient, a $\pm0.5\%$
velocity gradient changes $G_9(i\!\leftarrow\! j)$ by $10\%$, and a realistic
$67\%$ gradient by $70$--$340\%$ across frequency and slowness.

Coarsening the reference mid-solve is therefore \emph{not} a small perturbation
of $\mathcal{P}^{z}$.  Whatever multiresolution scheme is adopted, that
measurement sets the tolerance it must respect --- and it is another reason to
prefer nested iteration, where the operator is fixed within each solve, over
scale refinement inside a single Krylov space.

% =============================================================================
\section{Recommendation}

\begin{enumerate}
\item Place the scale hierarchy \emph{between} solves (nested iteration /
      cascadic), not within one Krylov space.  If it must go inside, use FGMRES
      or an explicit inexact-Krylov bound, and note that~\eqref{eq:relax} runs
      the opposite way to the intuition.
\item Coarsen the entries of $\bA$, never $(\alpha,\beta,\rho)$.  Haar is then
      the right tool and is second-order accurate.
\item For better than second order, use Backus/Capdeville homogenisation, and
      accept that the effective medium is anisotropic.
\item For operator compression rather than medium representation, Haar's single
      vanishing moment is inadequate for $1/r$ and $1/r^3$ kernels; use a basis
      with more vanishing moments in the Beylkin--Coifman--Rokhlin sense
      \citep{beylkin1991}.
\item Note that the two-potential formulation is \emph{already} a two-level
      scale split --- smooth reference plus fine contrast.  The natural
      generalisation is a hierarchy of such splits with coarse levels built by
      homogenisation.
\end{enumerate}

% =============================================================================
\begin{thebibliography}{99}

\bibitem[Backus(1962)]{backus1962}
Backus, G.~E. (1962).
Long-wave elastic anisotropy produced by horizontal layering.
\emph{Journal of Geophysical Research}, 67(11), 4427--4440.

\bibitem[Beylkin et al.(1991)]{beylkin1991}
Beylkin, G., Coifman, R., \& Rokhlin, V. (1991).
Fast wavelet transforms and numerical algorithms~I.
\emph{Communications on Pure and Applied Mathematics}, 44(2), 141--183.

\bibitem[Bornemann \& Deuflhard(1996)]{bornemann1996}
Bornemann, F.~A., \& Deuflhard, P. (1996).
The cascadic multigrid method for elliptic problems.
\emph{Numerische Mathematik}, 75(2), 135--152.

\bibitem[Bouras \& Frayss\'e(2005)]{bouras2005}
Bouras, A., \& Frayss\'e, V. (2005).
Inexact matrix--vector products in Krylov methods for solving linear systems:
a relaxation strategy.
\emph{SIAM Journal on Matrix Analysis and Applications}, 26(3), 660--678.
\url{https://epubs.siam.org/doi/10.1137/S0895479801384743}

\bibitem[Bunks et al.(1995)]{bunks1995}
Bunks, C., Saleck, F.~M., Zaleski, S., \& Chavent, G. (1995).
Multiscale seismic waveform inversion.
\emph{Geophysics}, 60(5), 1457--1473.

\bibitem[Capdeville et al.(2010)]{capdeville2010}
Capdeville, Y., Guillot, L., \& Marigo, J.-J. (2010).
2-D non-periodic homogenization to upscale elastic media for P--SV waves.
\emph{Geophysical Journal International}, 182(2), 903--922.
\url{https://academic.oup.com/gji/article/182/2/903/572295}

\bibitem[Capdeville et al.(2018)]{capdeville2018}
Capdeville, Y., Cupillard, P., \& Singh, S. (2018).
Non-periodic homogenization of 3-D elastic media for the seismic wave equation.
\emph{Geophysical Journal International}, 213(2), 983--1001.
\url{https://academic.oup.com/gji/article/213/2/983/4831486}

\bibitem[Daubechies(1992)]{daubechies1992}
Daubechies, I. (1992).
\emph{Ten Lectures on Wavelets}.
SIAM, Philadelphia.

\bibitem[Hustedt et al.(2004)]{hustedt2004}
Hustedt, B., Operto, S., \& Virieux, J. (2004).
A multi-level direct--iterative solver for seismic wave propagation modelling:
space and wavelet approaches.
\emph{Geophysical Journal International}, 155(3), 953--980.
\url{https://academic.oup.com/gji/article/155/3/953/632421}

\bibitem[Mallat(1989)]{mallat1989}
Mallat, S.~G. (1989).
A theory for multiresolution signal decomposition: the wavelet representation.
\emph{IEEE Transactions on Pattern Analysis and Machine Intelligence},
11(7), 674--693.

\bibitem[Saad(1993)]{saad1993}
Saad, Y. (1993).
A flexible inner--outer preconditioned GMRES algorithm.
\emph{SIAM Journal on Scientific Computing}, 14(2), 461--469.

\bibitem[Simoncini \& Szyld(2003)]{simoncini2003}
Simoncini, V., \& Szyld, D.~B. (2003).
Theory of inexact Krylov subspace methods and applications to scientific
computing.
\emph{SIAM Journal on Scientific Computing}, 25(2), 454--477.

\bibitem[van den Eshof \& Sleijpen(2004)]{vandeneshof2004}
van den Eshof, J., \& Sleijpen, G.~L.~G. (2004).
Inexact Krylov subspace methods for linear systems.
\emph{SIAM Journal on Matrix Analysis and Applications}, 26(1), 125--153.
\url{https://webspace.science.uu.nl/~sleij101/Reprints/ESl04b-40345.pdf}

\end{thebibliography}

\end{document}
